% !PRINTED none
% Abschnitt-opening page (part title + start of § 41); no running-head page
% number visible in any of the three bands supplied for this page.

\begin{center}
\gesperrt{Vierter Abschnitt.}

\vspace{1em}
{\Large\bfseries Die elliptischen Funktionen.}

\vspace{0.5em}
\rule{0.3\textwidth}{0.4pt}
\end{center}

\begin{center}
\textbf{\S\,41. \gesperrt{Zusammenhang der \(\vartheta\)-Funktionen mit den elliptischen\\
Integralen.}}
\end{center}

Nach \S\,21 bestehen zwischen den Quadraten der vier \(\vartheta\)-Funktionen
zwei voneinander unabhängige lineare Gleichungen, und man kann also zwei
von diesen Quadraten durch die beiden anderen oder auch alle vier durch
zwei unabhängige Variable \(\xi\), \(\eta\) ausdrücken. Indem wir das letztere tun,
bezeichnen wir mit \(\xi_1\), \(\eta_1\); \(\xi_2\), \(\eta_2\); \(\xi_3\), \(\eta_3\); \(\xi_4\), \(\eta_4\)
Konstanten und setzen, indem wir an die Bezeichnungsweise Bd.\,I, \S\,67
anknüpfen:
\weq{1}{\begin{aligned}
\vartheta^2_{01}(u) &= \xi\eta_1 - \eta\xi_1 = (\xi\eta_1),\\
\vartheta^2_{11}(u) &= \xi\eta_2 - \eta\xi_2 = (\xi\eta_2),\\
\vartheta^2_{10}(u) &= \xi\eta_3 - \eta\xi_3 = (\xi\eta_3),\\
\vartheta^2_{00}(u) &= \xi\eta_4 - \eta\xi_4 = (\xi\eta_4).
\end{aligned}}
Zwischen den Konstanten \(\xi_i\), \(\eta_i\) und den Werten \(\vartheta_{01}\), \(\vartheta_{10}\),
\(\vartheta_{00}\) bestehen vier Relationen, die sich aus den Gleichungen (13) des
\S\,21 herleiten lassen. Diese Gleichungen können wir nach der Bezeichnung
(1) in der Form schreiben:
\weq{2}{\begin{aligned}
(\xi\eta_3)\vartheta^2_{01} &= (\xi\eta_1)\vartheta^2_{10} - (\xi\eta_2)\vartheta^2_{00},\\
(\xi\eta_4)\vartheta^2_{01} &= (\xi\eta_1)\vartheta^2_{00} - (\xi\eta_2)\vartheta^2_{10}.
\end{aligned}}
Man kann diesen Relationen, indem man in (2) \(\xi\), \(\eta = \xi_1\), \(\eta_1\)
und \(= \xi_2\), \(\eta_2\) setzt, die Form geben:
\weq{3}{\begin{aligned}
(\xi_1\eta_3)\vartheta^2_{01} &= (\xi_2\eta_1)\vartheta^2_{00},\\
(\xi_2\eta_3)\vartheta^2_{01} &= (\xi_2\eta_1)\vartheta^2_{10},\\
(\xi_1\eta_4)\vartheta^2_{01} &= (\xi_2\eta_1)\vartheta^2_{10},\\
(\xi_2\eta_4)\vartheta^2_{01} &= (\xi_2\eta_1)\vartheta^2_{00},
\end{aligned}}
wozu man noch fügen kann, indem man in (2) \(\xi\), \(\eta = \xi_4\), \(\eta_4\) setzt
und (3) und \S\,21, (14) benutzt:
\weq{4}{(\xi_4\eta_3) = (\xi_2\eta_1).}
